
\documentclass{article}
\usepackage{colortbl}
\usepackage{makecell}
\usepackage{multirow}
\usepackage{supertabular}

\begin{document}

\newcounter{utterance}

\centering \large Interaction Transcript for game `codenames', experiment `opponent\_none', episode 1 with glm{-}5.2.
\vspace{24pt}

{ \footnotesize  \setcounter{utterance}{1}
\setlength{\tabcolsep}{0pt}
\begin{supertabular}{c@{$\;$}|p{.15\linewidth}@{}p{.15\linewidth}p{.15\linewidth}p{.15\linewidth}p{.15\linewidth}p{.15\linewidth}}
    \# & $\;$A & \multicolumn{4}{c}{Game Master} & $\;\:$B\\
    \hline

    \theutterance \stepcounter{utterance}  
    & & & \multicolumn{2}{p{0.3\linewidth}}{
        \cellcolor[rgb]{0.9,0.9,0.9}{
            \makecell[{{p{\linewidth}}}]{
                \texttt{\tiny{[GM$|$GM]}}
                {'hidden': {'team': ['птицечовка', 'антарктида', 'бомба', 'живот', 'пекин', 'оранжев', 'пистолет', 'ученик', 'мента'], 'innocent': ['вакуум', 'шапка', 'москва', 'промяна', 'карта', 'низ', 'поток'], 'opponent': ['ястреб', 'сила', 'пилот', 'триъгълник', 'двор', 'нож', 'лъч', 'връзка'], 'assassin': ['част']}, 'revealed': {'team': {'team': [], 'innocent': [], 'opponent': [], 'assassin': []}, 'opponent': {'team': [], 'innocent': [], 'opponent': [], 'assassin': []}}}
            }
        }
    }
    & & \\ \\

    \theutterance \stepcounter{utterance}  
    & & \multicolumn{4}{p{0.6\linewidth}}{
        \cellcolor[rgb]{0.9,0.9,0.9}{
            \makecell[{{p{\linewidth}}}]{
                \texttt{\tiny{[P1$\langle$GM]}}
                \texttt{Дайте една подсказка с една дума, която е свързана с някои от думите на вашия екип от списъка по{-}долу. Можете да избирате с колко и с кои думи да е свързана подсказката ви.} \\
\texttt{Целта е партньорът ви да отгатне всички думи на вашия отбор въз основа на подсказките, които му давате, преди другият отбор да е разкрил всички свои думи. Ако партньорът ви отгатне дума на "противника" или "разсейваща" дума, неговият ход за отгатване приключва и противниковият отбор получава възможност да направи ход. Ако партньорът ви отгатне дума "убиец", играта приключва и отборът ви губи незабавно.} \\
\texttt{Колкото повече думи свързва вашата подсказка, толкова по{-}бързо можете да спечелите играта, но асоциациите може да се окажат по{-}трудни за отгатване от партньора ви, така че избирайте подсказващата дума разумно. Подсказващата дума трябва да е семантично свързана с целевите думи, не може да е една от думите в списъците или да съдържа части от тях.} \\
\\ 
\texttt{Винаги давайте подсказката от една дума и списъка със свързаните с нея целеви думи, разделени със запетая, в следния формат и правете отговорите си възможно най{-}кратки, като никога не включвате друг текст, освен този, който се изисква в тази форма:} \\
\\ 
\texttt{ПОДСКАЗКА: <ДУМА>} \\
\texttt{ЦЕЛИ: <ЦЕЛИ>} \\
\\ 
\texttt{Думите на вашия отбор са: птицечовка, антарктида, бомба, живот, пекин, оранжев, пистолет, ученик, мента.} \\
\texttt{Думите на отбора на опонента ви са: ястреб, сила, пилот, триъгълник, двор, нож, лъч, връзка.} \\
\texttt{Разсейващите думи са: вакуум, шапка, москва, промяна, карта, низ, поток.} \\
\texttt{Думите убийци са: част.} \\
            }
        }
    }
    & \\ \\

    \theutterance \stepcounter{utterance}  
    & \multicolumn{4}{p{0.6\linewidth}}{
        \cellcolor[rgb]{0.9,0.9,0.9}{
            \makecell[{{p{\linewidth}}}]{
                \texttt{\tiny{[P1$\rangle$GM]}}
                \texttt{ПОДСКАЗКА: китай} \\
\texttt{ЦЕЛИ: пекин, птицечовка} \\
            }
        }
    }
    & & \\ \\

    \theutterance \stepcounter{utterance}  
    & & & \multicolumn{2}{p{0.3\linewidth}}{
        \cellcolor[rgb]{0.9,0.9,0.9}{
            \makecell[{{p{\linewidth}}}]{
                \texttt{\tiny{[GM$|$GM]}}
                \texttt{китай} \\
            }
        }
    }
    & & \\ \\

    \theutterance \stepcounter{utterance}  
    & & & \multicolumn{2}{p{0.3\linewidth}}{
        \cellcolor[rgb]{0.9,0.9,0.9}{
            \makecell[{{p{\linewidth}}}]{
                \texttt{\tiny{[GM$|$GM]}}
                ['пекин', 'птицечовка']
            }
        }
    }
    & & \\ \\

    \theutterance \stepcounter{utterance}  
    & & & \multicolumn{2}{p{0.3\linewidth}}{
        \cellcolor[rgb]{0.9,0.9,0.9}{
            \makecell[{{p{\linewidth}}}]{
                \texttt{\tiny{[GM$|$GM]}}
                {'word': 'пекин', 'assignment': 'team'}
            }
        }
    }
    & & \\ \\

    \theutterance \stepcounter{utterance}  
    & & & \multicolumn{2}{p{0.3\linewidth}}{
        \cellcolor[rgb]{0.9,0.9,0.9}{
            \makecell[{{p{\linewidth}}}]{
                \texttt{\tiny{[GM$|$GM]}}
                {'word': 'птицечовка', 'assignment': 'team'}
            }
        }
    }
    & & \\ \\

    \theutterance \stepcounter{utterance}  
    & & \multicolumn{4}{p{0.6\linewidth}}{
        \cellcolor[rgb]{0.9,0.9,0.9}{
            \makecell[{{p{\linewidth}}}]{
                \texttt{\tiny{[GM$\rangle$P2]}}
                \texttt{Представете списък, разделен със запетаи, от до 2 думи от следния списък, които най{-}добре се отнасят или са най{-}тясно свързани с думата 'китай'. Винаги започвайте списъка си с предположения с ОТГОВОР:  и не включвайте друг текст в отговора си.} \\
\\ 
\texttt{сила, антарктида, ястреб, низ, ученик, птицечовка, оранжев, бомба, промяна, пилот, връзка, москва, нож, шапка, живот, двор, триъгълник, пекин, вакуум, част, пистолет, карта, поток, мента, лъч} \\
            }
        }
    }
    & \\ \\

    \theutterance \stepcounter{utterance}  
    & & & \multicolumn{4}{p{0.6\linewidth}}{
        \cellcolor[rgb]{0.9,0.9,0.9}{
            \makecell[{{p{\linewidth}}}]{
                \texttt{\tiny{[GM$\langle$P2]}}
                \texttt{ОТГОВОР: пекин, сила} \\
            }
        }
    }
     \\ \\

    \theutterance \stepcounter{utterance}  
    & & & \multicolumn{2}{p{0.3\linewidth}}{
        \cellcolor[rgb]{0.9,0.9,0.9}{
            \makecell[{{p{\linewidth}}}]{
                \texttt{\tiny{[GM$|$GM]}}
                ['пекин', 'сила']
            }
        }
    }
    & & \\ \\

    \theutterance \stepcounter{utterance}  
    & & & \multicolumn{2}{p{0.3\linewidth}}{
        \cellcolor[rgb]{0.9,0.9,0.9}{
            \makecell[{{p{\linewidth}}}]{
                \texttt{\tiny{[GM$|$GM]}}
                {'word': 'пекин', 'assignment': 'team'}
            }
        }
    }
    & & \\ \\

    \theutterance \stepcounter{utterance}  
    & & & \multicolumn{2}{p{0.3\linewidth}}{
        \cellcolor[rgb]{0.9,0.9,0.9}{
            \makecell[{{p{\linewidth}}}]{
                \texttt{\tiny{[GM$|$GM]}}
                {'word': 'пекин', 'assignment': 'team'}
            }
        }
    }
    & & \\ \\

    \theutterance \stepcounter{utterance}  
    & & & \multicolumn{2}{p{0.3\linewidth}}{
        \cellcolor[rgb]{0.9,0.9,0.9}{
            \makecell[{{p{\linewidth}}}]{
                \texttt{\tiny{[GM$|$GM]}}
                {'word': 'сила', 'assignment': 'opponent'}
            }
        }
    }
    & & \\ \\

    \theutterance \stepcounter{utterance}  
    & & & \multicolumn{2}{p{0.3\linewidth}}{
        \cellcolor[rgb]{0.9,0.9,0.9}{
            \makecell[{{p{\linewidth}}}]{
                \texttt{\tiny{[GM$|$GM]}}
                \texttt{сила} \\
            }
        }
    }
    & & \\ \\

    \theutterance \stepcounter{utterance}  
    & & & \multicolumn{2}{p{0.3\linewidth}}{
        \cellcolor[rgb]{0.9,0.9,0.9}{
            \makecell[{{p{\linewidth}}}]{
                \texttt{\tiny{[GM$|$GM]}}
                {'hidden': {'team': ['птицечовка', 'антарктида', 'бомба', 'живот', 'оранжев', 'пистолет', 'ученик', 'мента'], 'innocent': ['вакуум', 'шапка', 'москва', 'промяна', 'карта', 'низ', 'поток'], 'opponent': ['ястреб', 'пилот', 'триъгълник', 'двор', 'нож', 'лъч', 'връзка'], 'assassin': ['част']}, 'revealed': {'team': {'team': ['пекин'], 'innocent': [], 'opponent': ['сила'], 'assassin': []}, 'opponent': {'team': [], 'innocent': [], 'opponent': [], 'assassin': []}}}
            }
        }
    }
    & & \\ \\

    \theutterance \stepcounter{utterance}  
    & & \multicolumn{4}{p{0.6\linewidth}}{
        \cellcolor[rgb]{0.9,0.9,0.9}{
            \makecell[{{p{\linewidth}}}]{
                \texttt{\tiny{[P1$\langle$GM]}}
                \texttt{Думата пекин беше позната правилно. Думата сила беше позната, но е противнически дума. Ходът на вашия отборник завърши там.} \\
\texttt{Сега дайте още една подсказка, свързана с някои от останалите думи на вашия отбор, и списък на свързаните целеви думи. Не забравяйте да започнете подсказката си с 'ПОДСКАЗКА: ', поставете нов ред и започнете списъка с целевите думи, разделен със запетая, с 'ЦЕЛИ: '. Забележете: някои думи са премахнати от списъците в сравнение с предишните заявки.} \\
\\ 
\texttt{Останалите думи на вашия отбор са: птицечовка, антарктида, бомба, живот, оранжев, пистолет, ученик, мента.} \\
\texttt{Останалите думи на вашия противник са: ястреб, пилот, триъгълник, двор, нож, лъч, връзка.} \\
\texttt{Останалите отвличащи думи са: вакуум, шапка, москва, промяна, карта, низ, поток.} \\
\texttt{Останалите убийствени думи са: част.} \\
            }
        }
    }
    & \\ \\

    \theutterance \stepcounter{utterance}  
    & \multicolumn{4}{p{0.6\linewidth}}{
        \cellcolor[rgb]{0.9,0.9,0.9}{
            \makecell[{{p{\linewidth}}}]{
                \texttt{\tiny{[P1$\rangle$GM]}}
                \texttt{ПОДСКАЗКА: оръжие} \\
\texttt{ЦЕЛИ: бомба, пистолет} \\
            }
        }
    }
    & & \\ \\

    \theutterance \stepcounter{utterance}  
    & & & \multicolumn{2}{p{0.3\linewidth}}{
        \cellcolor[rgb]{0.9,0.9,0.9}{
            \makecell[{{p{\linewidth}}}]{
                \texttt{\tiny{[GM$|$GM]}}
                \texttt{оръжие} \\
            }
        }
    }
    & & \\ \\

    \theutterance \stepcounter{utterance}  
    & & & \multicolumn{2}{p{0.3\linewidth}}{
        \cellcolor[rgb]{0.9,0.9,0.9}{
            \makecell[{{p{\linewidth}}}]{
                \texttt{\tiny{[GM$|$GM]}}
                ['бомба', 'пистолет']
            }
        }
    }
    & & \\ \\

    \theutterance \stepcounter{utterance}  
    & & & \multicolumn{2}{p{0.3\linewidth}}{
        \cellcolor[rgb]{0.9,0.9,0.9}{
            \makecell[{{p{\linewidth}}}]{
                \texttt{\tiny{[GM$|$GM]}}
                {'word': 'бомба', 'assignment': 'team'}
            }
        }
    }
    & & \\ \\

    \theutterance \stepcounter{utterance}  
    & & & \multicolumn{2}{p{0.3\linewidth}}{
        \cellcolor[rgb]{0.9,0.9,0.9}{
            \makecell[{{p{\linewidth}}}]{
                \texttt{\tiny{[GM$|$GM]}}
                {'word': 'пистолет', 'assignment': 'team'}
            }
        }
    }
    & & \\ \\

    \theutterance \stepcounter{utterance}  
    & & \multicolumn{4}{p{0.6\linewidth}}{
        \cellcolor[rgb]{0.9,0.9,0.9}{
            \makecell[{{p{\linewidth}}}]{
                \texttt{\tiny{[GM$\rangle$P2]}}
                \texttt{Думата пекин беше позната правилно. Думата сила беше позната, но е противнически дума. Вашият ход завърши там.} \\
\texttt{Сега предоставете още един списък, разделен със запетаи, от най{-}малко 1 и до 2 думи от следния списък с думи, които най{-}добре се отнасят или са най{-}тясно свързани с думата 'оръжие'. Не забравяйте да започнете отговора си с ОТГОВОР: . Забележете: някои думи са премахнати от списъка в сравнение с предишните заявки.} \\
\\ 
\texttt{антарктида, ястреб, низ, ученик, птицечовка, оранжев, бомба, промяна, пилот, връзка, москва, нож, шапка, живот, двор, триъгълник, вакуум, част, пистолет, карта, поток, мента, лъч} \\
            }
        }
    }
    & \\ \\

    \theutterance \stepcounter{utterance}  
    & & & \multicolumn{4}{p{0.6\linewidth}}{
        \cellcolor[rgb]{0.9,0.9,0.9}{
            \makecell[{{p{\linewidth}}}]{
                \texttt{\tiny{[GM$\langle$P2]}}
                \texttt{ОТГОВОР: пистолет, бомба} \\
            }
        }
    }
     \\ \\

    \theutterance \stepcounter{utterance}  
    & & & \multicolumn{2}{p{0.3\linewidth}}{
        \cellcolor[rgb]{0.9,0.9,0.9}{
            \makecell[{{p{\linewidth}}}]{
                \texttt{\tiny{[GM$|$GM]}}
                ['пистолет', 'бомба']
            }
        }
    }
    & & \\ \\

    \theutterance \stepcounter{utterance}  
    & & & \multicolumn{2}{p{0.3\linewidth}}{
        \cellcolor[rgb]{0.9,0.9,0.9}{
            \makecell[{{p{\linewidth}}}]{
                \texttt{\tiny{[GM$|$GM]}}
                {'word': 'пистолет', 'assignment': 'team'}
            }
        }
    }
    & & \\ \\

    \theutterance \stepcounter{utterance}  
    & & & \multicolumn{2}{p{0.3\linewidth}}{
        \cellcolor[rgb]{0.9,0.9,0.9}{
            \makecell[{{p{\linewidth}}}]{
                \texttt{\tiny{[GM$|$GM]}}
                {'word': 'пистолет', 'assignment': 'team'}
            }
        }
    }
    & & \\ \\

    \theutterance \stepcounter{utterance}  
    & & & \multicolumn{2}{p{0.3\linewidth}}{
        \cellcolor[rgb]{0.9,0.9,0.9}{
            \makecell[{{p{\linewidth}}}]{
                \texttt{\tiny{[GM$|$GM]}}
                {'word': 'бомба', 'assignment': 'team'}
            }
        }
    }
    & & \\ \\

    \theutterance \stepcounter{utterance}  
    & & & \multicolumn{2}{p{0.3\linewidth}}{
        \cellcolor[rgb]{0.9,0.9,0.9}{
            \makecell[{{p{\linewidth}}}]{
                \texttt{\tiny{[GM$|$GM]}}
                {'word': 'бомба', 'assignment': 'team'}
            }
        }
    }
    & & \\ \\

    \theutterance \stepcounter{utterance}  
    & & & \multicolumn{2}{p{0.3\linewidth}}{
        \cellcolor[rgb]{0.9,0.9,0.9}{
            \makecell[{{p{\linewidth}}}]{
                \texttt{\tiny{[GM$|$GM]}}
                {'hidden': {'team': ['птицечовка', 'антарктида', 'живот', 'оранжев', 'ученик', 'мента'], 'innocent': ['вакуум', 'шапка', 'москва', 'промяна', 'карта', 'низ', 'поток'], 'opponent': ['ястреб', 'пилот', 'триъгълник', 'двор', 'нож', 'лъч', 'връзка'], 'assassin': ['част']}, 'revealed': {'team': {'team': ['пекин', 'пистолет', 'бомба'], 'innocent': [], 'opponent': ['сила'], 'assassin': []}, 'opponent': {'team': [], 'innocent': [], 'opponent': [], 'assassin': []}}}
            }
        }
    }
    & & \\ \\

    \theutterance \stepcounter{utterance}  
    & & \multicolumn{4}{p{0.6\linewidth}}{
        \cellcolor[rgb]{0.9,0.9,0.9}{
            \makecell[{{p{\linewidth}}}]{
                \texttt{\tiny{[P1$\langle$GM]}}
                \texttt{Думите пистолет, бомба бяха познати правилно. Ходът на вашия отборник завърши там.} \\
\texttt{Сега дайте още една подсказка, свързана с някои от останалите думи на вашия отбор, и списък на свързаните целеви думи. Не забравяйте да започнете подсказката си с 'ПОДСКАЗКА: ', поставете нов ред и започнете списъка с целевите думи, разделен със запетая, с 'ЦЕЛИ: '. Забележете: някои думи са премахнати от списъците в сравнение с предишните заявки.} \\
\\ 
\texttt{Останалите думи на вашия отбор са: птицечовка, антарктида, живот, оранжев, ученик, мента.} \\
\texttt{Останалите думи на вашия противник са: ястреб, пилот, триъгълник, двор, нож, лъч, връзка.} \\
\texttt{Останалите отвличащи думи са: вакуум, шапка, москва, промяна, карта, низ, поток.} \\
\texttt{Останалите убийствени думи са: част.} \\
            }
        }
    }
    & \\ \\

    \theutterance \stepcounter{utterance}  
    & \multicolumn{4}{p{0.6\linewidth}}{
        \cellcolor[rgb]{0.9,0.9,0.9}{
            \makecell[{{p{\linewidth}}}]{
                \texttt{\tiny{[P1$\rangle$GM]}}
                \texttt{ПОДСКАЗКА: експедиция} \\
\texttt{ЦЕЛИ: антарктида, ученик} \\
            }
        }
    }
    & & \\ \\

    \theutterance \stepcounter{utterance}  
    & & & \multicolumn{2}{p{0.3\linewidth}}{
        \cellcolor[rgb]{0.9,0.9,0.9}{
            \makecell[{{p{\linewidth}}}]{
                \texttt{\tiny{[GM$|$GM]}}
                \texttt{експедиция} \\
            }
        }
    }
    & & \\ \\

    \theutterance \stepcounter{utterance}  
    & & & \multicolumn{2}{p{0.3\linewidth}}{
        \cellcolor[rgb]{0.9,0.9,0.9}{
            \makecell[{{p{\linewidth}}}]{
                \texttt{\tiny{[GM$|$GM]}}
                ['антарктида', 'ученик']
            }
        }
    }
    & & \\ \\

    \theutterance \stepcounter{utterance}  
    & & & \multicolumn{2}{p{0.3\linewidth}}{
        \cellcolor[rgb]{0.9,0.9,0.9}{
            \makecell[{{p{\linewidth}}}]{
                \texttt{\tiny{[GM$|$GM]}}
                {'word': 'антарктида', 'assignment': 'team'}
            }
        }
    }
    & & \\ \\

    \theutterance \stepcounter{utterance}  
    & & & \multicolumn{2}{p{0.3\linewidth}}{
        \cellcolor[rgb]{0.9,0.9,0.9}{
            \makecell[{{p{\linewidth}}}]{
                \texttt{\tiny{[GM$|$GM]}}
                {'word': 'ученик', 'assignment': 'team'}
            }
        }
    }
    & & \\ \\

    \theutterance \stepcounter{utterance}  
    & & \multicolumn{4}{p{0.6\linewidth}}{
        \cellcolor[rgb]{0.9,0.9,0.9}{
            \makecell[{{p{\linewidth}}}]{
                \texttt{\tiny{[GM$\rangle$P2]}}
                \texttt{Думите пистолет, бомба бяха познати правилно. Вашият ход завърши там.} \\
\texttt{Сега предоставете още един списък, разделен със запетаи, от най{-}малко 1 и до 2 думи от следния списък с думи, които най{-}добре се отнасят или са най{-}тясно свързани с думата 'експедиция'. Не забравяйте да започнете отговора си с ОТГОВОР: . Забележете: някои думи са премахнати от списъка в сравнение с предишните заявки.} \\
\\ 
\texttt{антарктида, ястреб, низ, ученик, птицечовка, оранжев, промяна, пилот, връзка, москва, нож, шапка, живот, двор, триъгълник, вакуум, част, карта, поток, мента, лъч} \\
            }
        }
    }
    & \\ \\

    \theutterance \stepcounter{utterance}  
    & & & \multicolumn{4}{p{0.6\linewidth}}{
        \cellcolor[rgb]{0.9,0.9,0.9}{
            \makecell[{{p{\linewidth}}}]{
                \texttt{\tiny{[GM$\langle$P2]}}
                \texttt{ОТГОВОР: антарктида, карта} \\
            }
        }
    }
     \\ \\

    \theutterance \stepcounter{utterance}  
    & & & \multicolumn{2}{p{0.3\linewidth}}{
        \cellcolor[rgb]{0.9,0.9,0.9}{
            \makecell[{{p{\linewidth}}}]{
                \texttt{\tiny{[GM$|$GM]}}
                ['антарктида', 'карта']
            }
        }
    }
    & & \\ \\

    \theutterance \stepcounter{utterance}  
    & & & \multicolumn{2}{p{0.3\linewidth}}{
        \cellcolor[rgb]{0.9,0.9,0.9}{
            \makecell[{{p{\linewidth}}}]{
                \texttt{\tiny{[GM$|$GM]}}
                {'word': 'антарктида', 'assignment': 'team'}
            }
        }
    }
    & & \\ \\

    \theutterance \stepcounter{utterance}  
    & & & \multicolumn{2}{p{0.3\linewidth}}{
        \cellcolor[rgb]{0.9,0.9,0.9}{
            \makecell[{{p{\linewidth}}}]{
                \texttt{\tiny{[GM$|$GM]}}
                {'word': 'антарктида', 'assignment': 'team'}
            }
        }
    }
    & & \\ \\

    \theutterance \stepcounter{utterance}  
    & & & \multicolumn{2}{p{0.3\linewidth}}{
        \cellcolor[rgb]{0.9,0.9,0.9}{
            \makecell[{{p{\linewidth}}}]{
                \texttt{\tiny{[GM$|$GM]}}
                {'word': 'карта', 'assignment': 'innocent'}
            }
        }
    }
    & & \\ \\

    \theutterance \stepcounter{utterance}  
    & & & \multicolumn{2}{p{0.3\linewidth}}{
        \cellcolor[rgb]{0.9,0.9,0.9}{
            \makecell[{{p{\linewidth}}}]{
                \texttt{\tiny{[GM$|$GM]}}
                \texttt{карта} \\
            }
        }
    }
    & & \\ \\

    \theutterance \stepcounter{utterance}  
    & & & \multicolumn{2}{p{0.3\linewidth}}{
        \cellcolor[rgb]{0.9,0.9,0.9}{
            \makecell[{{p{\linewidth}}}]{
                \texttt{\tiny{[GM$|$GM]}}
                {'hidden': {'team': ['птицечовка', 'живот', 'оранжев', 'ученик', 'мента'], 'innocent': ['вакуум', 'шапка', 'москва', 'промяна', 'низ', 'поток'], 'opponent': ['ястреб', 'пилот', 'триъгълник', 'двор', 'нож', 'лъч', 'връзка'], 'assassin': ['част']}, 'revealed': {'team': {'team': ['пекин', 'пистолет', 'бомба', 'антарктида'], 'innocent': ['карта'], 'opponent': ['сила'], 'assassin': []}, 'opponent': {'team': [], 'innocent': [], 'opponent': [], 'assassin': []}}}
            }
        }
    }
    & & \\ \\

    \theutterance \stepcounter{utterance}  
    & & \multicolumn{4}{p{0.6\linewidth}}{
        \cellcolor[rgb]{0.9,0.9,0.9}{
            \makecell[{{p{\linewidth}}}]{
                \texttt{\tiny{[P1$\langle$GM]}}
                \texttt{Думата антарктида беше позната правилно. Думата карта беше позната, но е неутрален дума. Ходът на вашия отборник завърши там.} \\
\texttt{Сега дайте още една подсказка, свързана с някои от останалите думи на вашия отбор, и списък на свързаните целеви думи. Не забравяйте да започнете подсказката си с 'ПОДСКАЗКА: ', поставете нов ред и започнете списъка с целевите думи, разделен със запетая, с 'ЦЕЛИ: '. Забележете: някои думи са премахнати от списъците в сравнение с предишните заявки.} \\
\\ 
\texttt{Останалите думи на вашия отбор са: птицечовка, живот, оранжев, ученик, мента.} \\
\texttt{Останалите думи на вашия противник са: ястреб, пилот, триъгълник, двор, нож, лъч, връзка.} \\
\texttt{Останалите отвличащи думи са: вакуум, шапка, москва, промяна, низ, поток.} \\
\texttt{Останалите убийствени думи са: част.} \\
            }
        }
    }
    & \\ \\

    \theutterance \stepcounter{utterance}  
    & \multicolumn{4}{p{0.6\linewidth}}{
        \cellcolor[rgb]{0.9,0.9,0.9}{
            \makecell[{{p{\linewidth}}}]{
                \texttt{\tiny{[P1$\rangle$GM]}}
                \texttt{ПОДСКАЗКА: школа} \\
\texttt{ЦЕЛИ: ученик, мента} \\
            }
        }
    }
    & & \\ \\

    \theutterance \stepcounter{utterance}  
    & & & \multicolumn{2}{p{0.3\linewidth}}{
        \cellcolor[rgb]{0.9,0.9,0.9}{
            \makecell[{{p{\linewidth}}}]{
                \texttt{\tiny{[GM$|$GM]}}
                \texttt{школа} \\
            }
        }
    }
    & & \\ \\

    \theutterance \stepcounter{utterance}  
    & & & \multicolumn{2}{p{0.3\linewidth}}{
        \cellcolor[rgb]{0.9,0.9,0.9}{
            \makecell[{{p{\linewidth}}}]{
                \texttt{\tiny{[GM$|$GM]}}
                ['ученик', 'мента']
            }
        }
    }
    & & \\ \\

    \theutterance \stepcounter{utterance}  
    & & & \multicolumn{2}{p{0.3\linewidth}}{
        \cellcolor[rgb]{0.9,0.9,0.9}{
            \makecell[{{p{\linewidth}}}]{
                \texttt{\tiny{[GM$|$GM]}}
                {'word': 'ученик', 'assignment': 'team'}
            }
        }
    }
    & & \\ \\

    \theutterance \stepcounter{utterance}  
    & & & \multicolumn{2}{p{0.3\linewidth}}{
        \cellcolor[rgb]{0.9,0.9,0.9}{
            \makecell[{{p{\linewidth}}}]{
                \texttt{\tiny{[GM$|$GM]}}
                {'word': 'мента', 'assignment': 'team'}
            }
        }
    }
    & & \\ \\

    \theutterance \stepcounter{utterance}  
    & & \multicolumn{4}{p{0.6\linewidth}}{
        \cellcolor[rgb]{0.9,0.9,0.9}{
            \makecell[{{p{\linewidth}}}]{
                \texttt{\tiny{[GM$\rangle$P2]}}
                \texttt{Думата антарктида беше позната правилно. Думата карта беше позната, но е неутрален дума. Вашият ход завърши там.} \\
\texttt{Сега предоставете още един списък, разделен със запетаи, от най{-}малко 1 и до 2 думи от следния списък с думи, които най{-}добре се отнасят или са най{-}тясно свързани с думата 'школа'. Не забравяйте да започнете отговора си с ОТГОВОР: . Забележете: някои думи са премахнати от списъка в сравнение с предишните заявки.} \\
\\ 
\texttt{ястреб, низ, ученик, птицечовка, оранжев, промяна, пилот, връзка, москва, нож, шапка, живот, двор, триъгълник, вакуум, част, поток, мента, лъч} \\
            }
        }
    }
    & \\ \\

    \theutterance \stepcounter{utterance}  
    & & & \multicolumn{4}{p{0.6\linewidth}}{
        \cellcolor[rgb]{0.9,0.9,0.9}{
            \makecell[{{p{\linewidth}}}]{
                \texttt{\tiny{[GM$\langle$P2]}}
                \texttt{ОТГОВОР: ученик, двор} \\
            }
        }
    }
     \\ \\

    \theutterance \stepcounter{utterance}  
    & & & \multicolumn{2}{p{0.3\linewidth}}{
        \cellcolor[rgb]{0.9,0.9,0.9}{
            \makecell[{{p{\linewidth}}}]{
                \texttt{\tiny{[GM$|$GM]}}
                ['ученик', 'двор']
            }
        }
    }
    & & \\ \\

    \theutterance \stepcounter{utterance}  
    & & & \multicolumn{2}{p{0.3\linewidth}}{
        \cellcolor[rgb]{0.9,0.9,0.9}{
            \makecell[{{p{\linewidth}}}]{
                \texttt{\tiny{[GM$|$GM]}}
                {'word': 'ученик', 'assignment': 'team'}
            }
        }
    }
    & & \\ \\

    \theutterance \stepcounter{utterance}  
    & & & \multicolumn{2}{p{0.3\linewidth}}{
        \cellcolor[rgb]{0.9,0.9,0.9}{
            \makecell[{{p{\linewidth}}}]{
                \texttt{\tiny{[GM$|$GM]}}
                {'word': 'ученик', 'assignment': 'team'}
            }
        }
    }
    & & \\ \\

    \theutterance \stepcounter{utterance}  
    & & & \multicolumn{2}{p{0.3\linewidth}}{
        \cellcolor[rgb]{0.9,0.9,0.9}{
            \makecell[{{p{\linewidth}}}]{
                \texttt{\tiny{[GM$|$GM]}}
                {'word': 'двор', 'assignment': 'opponent'}
            }
        }
    }
    & & \\ \\

    \theutterance \stepcounter{utterance}  
    & & & \multicolumn{2}{p{0.3\linewidth}}{
        \cellcolor[rgb]{0.9,0.9,0.9}{
            \makecell[{{p{\linewidth}}}]{
                \texttt{\tiny{[GM$|$GM]}}
                \texttt{двор} \\
            }
        }
    }
    & & \\ \\

    \theutterance \stepcounter{utterance}  
    & & & \multicolumn{2}{p{0.3\linewidth}}{
        \cellcolor[rgb]{0.9,0.9,0.9}{
            \makecell[{{p{\linewidth}}}]{
                \texttt{\tiny{[GM$|$GM]}}
                {'hidden': {'team': ['птицечовка', 'живот', 'оранжев', 'мента'], 'innocent': ['вакуум', 'шапка', 'москва', 'промяна', 'низ', 'поток'], 'opponent': ['ястреб', 'пилот', 'триъгълник', 'нож', 'лъч', 'връзка'], 'assassin': ['част']}, 'revealed': {'team': {'team': ['пекин', 'пистолет', 'бомба', 'антарктида', 'ученик'], 'innocent': ['карта'], 'opponent': ['сила', 'двор'], 'assassin': []}, 'opponent': {'team': [], 'innocent': [], 'opponent': [], 'assassin': []}}}
            }
        }
    }
    & & \\ \\

    \theutterance \stepcounter{utterance}  
    & & \multicolumn{4}{p{0.6\linewidth}}{
        \cellcolor[rgb]{0.9,0.9,0.9}{
            \makecell[{{p{\linewidth}}}]{
                \texttt{\tiny{[P1$\langle$GM]}}
                \texttt{Думата ученик беше позната правилно. Думата двор беше позната, но е противнически дума. Ходът на вашия отборник завърши там.} \\
\texttt{Сега дайте още една подсказка, свързана с някои от останалите думи на вашия отбор, и списък на свързаните целеви думи. Не забравяйте да започнете подсказката си с 'ПОДСКАЗКА: ', поставете нов ред и започнете списъка с целевите думи, разделен със запетая, с 'ЦЕЛИ: '. Забележете: някои думи са премахнати от списъците в сравнение с предишните заявки.} \\
\\ 
\texttt{Останалите думи на вашия отбор са: птицечовка, живот, оранжев, мента.} \\
\texttt{Останалите думи на вашия противник са: ястреб, пилот, триъгълник, нож, лъч, връзка.} \\
\texttt{Останалите отвличащи думи са: вакуум, шапка, москва, промяна, низ, поток.} \\
\texttt{Останалите убийствени думи са: част.} \\
            }
        }
    }
    & \\ \\

    \theutterance \stepcounter{utterance}  
    & \multicolumn{4}{p{0.6\linewidth}}{
        \cellcolor[rgb]{0.9,0.9,0.9}{
            \makecell[{{p{\linewidth}}}]{
                \texttt{\tiny{[P1$\rangle$GM]}}
                \texttt{ПОДСКАЗКА: плод} \\
\texttt{ЦЕЛИ: оранжев, мента} \\
            }
        }
    }
    & & \\ \\

    \theutterance \stepcounter{utterance}  
    & & & \multicolumn{2}{p{0.3\linewidth}}{
        \cellcolor[rgb]{0.9,0.9,0.9}{
            \makecell[{{p{\linewidth}}}]{
                \texttt{\tiny{[GM$|$GM]}}
                \texttt{плод} \\
            }
        }
    }
    & & \\ \\

    \theutterance \stepcounter{utterance}  
    & & & \multicolumn{2}{p{0.3\linewidth}}{
        \cellcolor[rgb]{0.9,0.9,0.9}{
            \makecell[{{p{\linewidth}}}]{
                \texttt{\tiny{[GM$|$GM]}}
                ['оранжев', 'мента']
            }
        }
    }
    & & \\ \\

    \theutterance \stepcounter{utterance}  
    & & & \multicolumn{2}{p{0.3\linewidth}}{
        \cellcolor[rgb]{0.9,0.9,0.9}{
            \makecell[{{p{\linewidth}}}]{
                \texttt{\tiny{[GM$|$GM]}}
                {'word': 'оранжев', 'assignment': 'team'}
            }
        }
    }
    & & \\ \\

    \theutterance \stepcounter{utterance}  
    & & & \multicolumn{2}{p{0.3\linewidth}}{
        \cellcolor[rgb]{0.9,0.9,0.9}{
            \makecell[{{p{\linewidth}}}]{
                \texttt{\tiny{[GM$|$GM]}}
                {'word': 'мента', 'assignment': 'team'}
            }
        }
    }
    & & \\ \\

    \theutterance \stepcounter{utterance}  
    & & \multicolumn{4}{p{0.6\linewidth}}{
        \cellcolor[rgb]{0.9,0.9,0.9}{
            \makecell[{{p{\linewidth}}}]{
                \texttt{\tiny{[GM$\rangle$P2]}}
                \texttt{Думата ученик беше позната правилно. Думата двор беше позната, но е противнически дума. Вашият ход завърши там.} \\
\texttt{Сега предоставете още един списък, разделен със запетаи, от най{-}малко 1 и до 2 думи от следния списък с думи, които най{-}добре се отнасят или са най{-}тясно свързани с думата 'плод'. Не забравяйте да започнете отговора си с ОТГОВОР: . Забележете: някои думи са премахнати от списъка в сравнение с предишните заявки.} \\
\\ 
\texttt{ястреб, низ, птицечовка, оранжев, промяна, пилот, връзка, москва, нож, шапка, живот, триъгълник, вакуум, част, поток, мента, лъч} \\
            }
        }
    }
    & \\ \\

    \theutterance \stepcounter{utterance}  
    & & & \multicolumn{4}{p{0.6\linewidth}}{
        \cellcolor[rgb]{0.9,0.9,0.9}{
            \makecell[{{p{\linewidth}}}]{
                \texttt{\tiny{[GM$\langle$P2]}}
                \texttt{ОТГОВОР: птицечовка, мента} \\
            }
        }
    }
     \\ \\

    \theutterance \stepcounter{utterance}  
    & & & \multicolumn{2}{p{0.3\linewidth}}{
        \cellcolor[rgb]{0.9,0.9,0.9}{
            \makecell[{{p{\linewidth}}}]{
                \texttt{\tiny{[GM$|$GM]}}
                ['птицечовка', 'мента']
            }
        }
    }
    & & \\ \\

    \theutterance \stepcounter{utterance}  
    & & & \multicolumn{2}{p{0.3\linewidth}}{
        \cellcolor[rgb]{0.9,0.9,0.9}{
            \makecell[{{p{\linewidth}}}]{
                \texttt{\tiny{[GM$|$GM]}}
                {'word': 'птицечовка', 'assignment': 'team'}
            }
        }
    }
    & & \\ \\

    \theutterance \stepcounter{utterance}  
    & & & \multicolumn{2}{p{0.3\linewidth}}{
        \cellcolor[rgb]{0.9,0.9,0.9}{
            \makecell[{{p{\linewidth}}}]{
                \texttt{\tiny{[GM$|$GM]}}
                {'word': 'мента', 'assignment': 'team'}
            }
        }
    }
    & & \\ \\

    \theutterance \stepcounter{utterance}  
    & & & \multicolumn{2}{p{0.3\linewidth}}{
        \cellcolor[rgb]{0.9,0.9,0.9}{
            \makecell[{{p{\linewidth}}}]{
                \texttt{\tiny{[GM$|$GM]}}
                {'word': 'мента', 'assignment': 'team'}
            }
        }
    }
    & & \\ \\

    \theutterance \stepcounter{utterance}  
    & & & \multicolumn{2}{p{0.3\linewidth}}{
        \cellcolor[rgb]{0.9,0.9,0.9}{
            \makecell[{{p{\linewidth}}}]{
                \texttt{\tiny{[GM$|$GM]}}
                {'hidden': {'team': ['живот', 'оранжев'], 'innocent': ['вакуум', 'шапка', 'москва', 'промяна', 'низ', 'поток'], 'opponent': ['ястреб', 'пилот', 'триъгълник', 'нож', 'лъч', 'връзка'], 'assassin': ['част']}, 'revealed': {'team': {'team': ['пекин', 'пистолет', 'бомба', 'антарктида', 'ученик', 'птицечовка', 'мента'], 'innocent': ['карта'], 'opponent': ['сила', 'двор'], 'assassin': []}, 'opponent': {'team': [], 'innocent': [], 'opponent': [], 'assassin': []}}}
            }
        }
    }
    & & \\ \\

    \theutterance \stepcounter{utterance}  
    & & \multicolumn{4}{p{0.6\linewidth}}{
        \cellcolor[rgb]{0.9,0.9,0.9}{
            \makecell[{{p{\linewidth}}}]{
                \texttt{\tiny{[P1$\langle$GM]}}
                \texttt{Думите птицечовка, мента бяха познати правилно. Ходът на вашия отборник завърши там.} \\
\texttt{Сега дайте още една подсказка, свързана с някои от останалите думи на вашия отбор, и списък на свързаните целеви думи. Не забравяйте да започнете подсказката си с 'ПОДСКАЗКА: ', поставете нов ред и започнете списъка с целевите думи, разделен със запетая, с 'ЦЕЛИ: '. Забележете: някои думи са премахнати от списъците в сравнение с предишните заявки.} \\
\\ 
\texttt{Останалите думи на вашия отбор са: живот, оранжев.} \\
\texttt{Останалите думи на вашия противник са: ястреб, пилот, триъгълник, нож, лъч, връзка.} \\
\texttt{Останалите отвличащи думи са: вакуум, шапка, москва, промяна, низ, поток.} \\
\texttt{Останалите убийствени думи са: част.} \\
            }
        }
    }
    & \\ \\

    \theutterance \stepcounter{utterance}  
    & \multicolumn{4}{p{0.6\linewidth}}{
        \cellcolor[rgb]{0.9,0.9,0.9}{
            \makecell[{{p{\linewidth}}}]{
                \texttt{\tiny{[P1$\rangle$GM]}}
                \texttt{ПОДСКАЗКА: цвят} \\
\texttt{ЦЕЛИ: живот, оранжев} \\
            }
        }
    }
    & & \\ \\

    \theutterance \stepcounter{utterance}  
    & & & \multicolumn{2}{p{0.3\linewidth}}{
        \cellcolor[rgb]{0.9,0.9,0.9}{
            \makecell[{{p{\linewidth}}}]{
                \texttt{\tiny{[GM$|$GM]}}
                \texttt{цвят} \\
            }
        }
    }
    & & \\ \\

    \theutterance \stepcounter{utterance}  
    & & & \multicolumn{2}{p{0.3\linewidth}}{
        \cellcolor[rgb]{0.9,0.9,0.9}{
            \makecell[{{p{\linewidth}}}]{
                \texttt{\tiny{[GM$|$GM]}}
                ['живот', 'оранжев']
            }
        }
    }
    & & \\ \\

    \theutterance \stepcounter{utterance}  
    & & & \multicolumn{2}{p{0.3\linewidth}}{
        \cellcolor[rgb]{0.9,0.9,0.9}{
            \makecell[{{p{\linewidth}}}]{
                \texttt{\tiny{[GM$|$GM]}}
                {'word': 'живот', 'assignment': 'team'}
            }
        }
    }
    & & \\ \\

    \theutterance \stepcounter{utterance}  
    & & & \multicolumn{2}{p{0.3\linewidth}}{
        \cellcolor[rgb]{0.9,0.9,0.9}{
            \makecell[{{p{\linewidth}}}]{
                \texttt{\tiny{[GM$|$GM]}}
                {'word': 'оранжев', 'assignment': 'team'}
            }
        }
    }
    & & \\ \\

    \theutterance \stepcounter{utterance}  
    & & \multicolumn{4}{p{0.6\linewidth}}{
        \cellcolor[rgb]{0.9,0.9,0.9}{
            \makecell[{{p{\linewidth}}}]{
                \texttt{\tiny{[GM$\rangle$P2]}}
                \texttt{Думите птицечовка, мента бяха познати правилно. Вашият ход завърши там.} \\
\texttt{Сега предоставете още един списък, разделен със запетаи, от най{-}малко 1 и до 2 думи от следния списък с думи, които най{-}добре се отнасят или са най{-}тясно свързани с думата 'цвят'. Не забравяйте да започнете отговора си с ОТГОВОР: . Забележете: някои думи са премахнати от списъка в сравнение с предишните заявки.} \\
\\ 
\texttt{ястреб, низ, оранжев, промяна, пилот, връзка, москва, нож, шапка, живот, триъгълник, вакуум, част, поток, лъч} \\
            }
        }
    }
    & \\ \\

    \theutterance \stepcounter{utterance}  
    & & & \multicolumn{4}{p{0.6\linewidth}}{
        \cellcolor[rgb]{0.9,0.9,0.9}{
            \makecell[{{p{\linewidth}}}]{
                \texttt{\tiny{[GM$\langle$P2]}}
                \texttt{ОТГОВОР: оранжев} \\
            }
        }
    }
     \\ \\

    \theutterance \stepcounter{utterance}  
    & & & \multicolumn{2}{p{0.3\linewidth}}{
        \cellcolor[rgb]{0.9,0.9,0.9}{
            \makecell[{{p{\linewidth}}}]{
                \texttt{\tiny{[GM$|$GM]}}
                ['оранжев']
            }
        }
    }
    & & \\ \\

    \theutterance \stepcounter{utterance}  
    & & & \multicolumn{2}{p{0.3\linewidth}}{
        \cellcolor[rgb]{0.9,0.9,0.9}{
            \makecell[{{p{\linewidth}}}]{
                \texttt{\tiny{[GM$|$GM]}}
                {'word': 'оранжев', 'assignment': 'team'}
            }
        }
    }
    & & \\ \\

    \theutterance \stepcounter{utterance}  
    & & & \multicolumn{2}{p{0.3\linewidth}}{
        \cellcolor[rgb]{0.9,0.9,0.9}{
            \makecell[{{p{\linewidth}}}]{
                \texttt{\tiny{[GM$|$GM]}}
                {'word': 'оранжев', 'assignment': 'team'}
            }
        }
    }
    & & \\ \\

    \theutterance \stepcounter{utterance}  
    & & & \multicolumn{2}{p{0.3\linewidth}}{
        \cellcolor[rgb]{0.9,0.9,0.9}{
            \makecell[{{p{\linewidth}}}]{
                \texttt{\tiny{[GM$|$GM]}}
                {'hidden': {'team': ['живот'], 'innocent': ['вакуум', 'шапка', 'москва', 'промяна', 'низ', 'поток'], 'opponent': ['ястреб', 'пилот', 'триъгълник', 'нож', 'лъч', 'връзка'], 'assassin': ['част']}, 'revealed': {'team': {'team': ['пекин', 'пистолет', 'бомба', 'антарктида', 'ученик', 'птицечовка', 'мента', 'оранжев'], 'innocent': ['карта'], 'opponent': ['сила', 'двор'], 'assassin': []}, 'opponent': {'team': [], 'innocent': [], 'opponent': [], 'assassin': []}}}
            }
        }
    }
    & & \\ \\

    \theutterance \stepcounter{utterance}  
    & & \multicolumn{4}{p{0.6\linewidth}}{
        \cellcolor[rgb]{0.9,0.9,0.9}{
            \makecell[{{p{\linewidth}}}]{
                \texttt{\tiny{[P1$\langle$GM]}}
                \texttt{Думата оранжев беше позната правилно. Ходът на вашия отборник завърши там.} \\
\texttt{Сега дайте още една подсказка, свързана с някои от останалите думи на вашия отбор, и списък на свързаните целеви думи. Не забравяйте да започнете подсказката си с 'ПОДСКАЗКА: ', поставете нов ред и започнете списъка с целевите думи, разделен със запетая, с 'ЦЕЛИ: '. Забележете: някои думи са премахнати от списъците в сравнение с предишните заявки.} \\
\\ 
\texttt{Останалите думи на вашия отбор са: живот.} \\
\texttt{Останалите думи на вашия противник са: ястреб, пилот, триъгълник, нож, лъч, връзка.} \\
\texttt{Останалите отвличащи думи са: вакуум, шапка, москва, промяна, низ, поток.} \\
\texttt{Останалите убийствени думи са: част.} \\
            }
        }
    }
    & \\ \\

    \theutterance \stepcounter{utterance}  
    & \multicolumn{4}{p{0.6\linewidth}}{
        \cellcolor[rgb]{0.9,0.9,0.9}{
            \makecell[{{p{\linewidth}}}]{
                \texttt{\tiny{[P1$\rangle$GM]}}
                \texttt{ПОДСКАЗКА: биология} \\
\texttt{ЦЕЛИ: живот} \\
            }
        }
    }
    & & \\ \\

    \theutterance \stepcounter{utterance}  
    & & & \multicolumn{2}{p{0.3\linewidth}}{
        \cellcolor[rgb]{0.9,0.9,0.9}{
            \makecell[{{p{\linewidth}}}]{
                \texttt{\tiny{[GM$|$GM]}}
                \texttt{биология} \\
            }
        }
    }
    & & \\ \\

    \theutterance \stepcounter{utterance}  
    & & & \multicolumn{2}{p{0.3\linewidth}}{
        \cellcolor[rgb]{0.9,0.9,0.9}{
            \makecell[{{p{\linewidth}}}]{
                \texttt{\tiny{[GM$|$GM]}}
                ['живот']
            }
        }
    }
    & & \\ \\

    \theutterance \stepcounter{utterance}  
    & & & \multicolumn{2}{p{0.3\linewidth}}{
        \cellcolor[rgb]{0.9,0.9,0.9}{
            \makecell[{{p{\linewidth}}}]{
                \texttt{\tiny{[GM$|$GM]}}
                {'word': 'живот', 'assignment': 'team'}
            }
        }
    }
    & & \\ \\

    \theutterance \stepcounter{utterance}  
    & & \multicolumn{4}{p{0.6\linewidth}}{
        \cellcolor[rgb]{0.9,0.9,0.9}{
            \makecell[{{p{\linewidth}}}]{
                \texttt{\tiny{[GM$\rangle$P2]}}
                \texttt{Думата оранжев беше позната правилно. Вашият ход завърши там.} \\
\texttt{Сега предоставете още един списък, разделен със запетаи, от най{-}малко 1 и до 1 думи от следния списък с думи, които най{-}добре се отнасят или са най{-}тясно свързани с думата 'биология'. Не забравяйте да започнете отговора си с ОТГОВОР: . Забележете: някои думи са премахнати от списъка в сравнение с предишните заявки.} \\
\\ 
\texttt{ястреб, низ, промяна, пилот, връзка, москва, нож, шапка, живот, триъгълник, вакуум, част, поток, лъч} \\
            }
        }
    }
    & \\ \\

    \theutterance \stepcounter{utterance}  
    & & & \multicolumn{4}{p{0.6\linewidth}}{
        \cellcolor[rgb]{0.9,0.9,0.9}{
            \makecell[{{p{\linewidth}}}]{
                \texttt{\tiny{[GM$\langle$P2]}}
                \texttt{ОТГОВОР: живот} \\
            }
        }
    }
     \\ \\

    \theutterance \stepcounter{utterance}  
    & & & \multicolumn{2}{p{0.3\linewidth}}{
        \cellcolor[rgb]{0.9,0.9,0.9}{
            \makecell[{{p{\linewidth}}}]{
                \texttt{\tiny{[GM$|$GM]}}
                ['живот']
            }
        }
    }
    & & \\ \\

    \theutterance \stepcounter{utterance}  
    & & & \multicolumn{2}{p{0.3\linewidth}}{
        \cellcolor[rgb]{0.9,0.9,0.9}{
            \makecell[{{p{\linewidth}}}]{
                \texttt{\tiny{[GM$|$GM]}}
                {'word': 'живот', 'assignment': 'team'}
            }
        }
    }
    & & \\ \\

    \theutterance \stepcounter{utterance}  
    & & & \multicolumn{2}{p{0.3\linewidth}}{
        \cellcolor[rgb]{0.9,0.9,0.9}{
            \makecell[{{p{\linewidth}}}]{
                \texttt{\tiny{[GM$|$GM]}}
                {'word': 'живот', 'assignment': 'team'}
            }
        }
    }
    & & \\ \\

    \theutterance \stepcounter{utterance}  
    & & & \multicolumn{2}{p{0.3\linewidth}}{
        \cellcolor[rgb]{0.9,0.9,0.9}{
            \makecell[{{p{\linewidth}}}]{
                \texttt{\tiny{[GM$|$GM]}}
                \texttt{team has won} \\
            }
        }
    }
    & & \\ \\

\end{supertabular}
}

\end{document}
